\documentclass[10pt,a4paper]{amsart}
\usepackage[margin=19mm]{geometry}
\usepackage{amsmath,amssymb,mathtools}
\usepackage[scr=rsfs,cal=euler]{mathalfa}
\usepackage{xfrac}
\usepackage[unicode,hidelinks,bookmarks=false]{hyperref}
\newcommand{\ie}{i.e.\ }
\newcommand{\cf}{cf.\ }
\newcommand{\resp}{respectively\ }
\newcommand{\Subsection}{\subsection}
\providecommand{\nobreakdash}{\nobreak\hbox{-}}
\setlength{\emergencystretch}{2em}
\multlinegap0pt
\usepackage{amssymb}
\usepackage{xcolor}
\newtheorem{proposition}{Proposition}
\newtheorem{lemma}{Lemma}
\title{The inverse curve shortening flow on the hyperbolic plane}
\author{Ivan Krznari\'c, Rafael L\'opez}
\date{Source: arXiv:2605.14385v1 (CC BY 4.0)}
\begin{document}
\maketitle

\noindent\emph{Source and licence.} The inverse curve shortening flow on the hyperbolic plane. Version arXiv:2605.14385v1 (CC BY 4.0), distributed by arXiv under the Creative Commons Attribution 4.0 International licence, http://creativecommons.org/licenses/by/4.0/, as recorded in the arXiv licence metadata for this exact version and shown on the abstract page https://arxiv.org/abs/2605.14385v1. Full licence notice: \texttt{licenses/arxiv-2605.14385-CC-BY-4.0.txt}.\par\smallskip

\noindent\textit{Editorial scope.} Counted unit: the article's lemma le56 (source0.tex lines 453--455). The article's complete original proof of it is delivered with it (source0.tex lines 457--477, one \texttt{proof} environment of 21 lines). No mathematics is written, repaired or re-derived: the statement and the proof are the article's own lines, in the article's own notation, and no retained line is substituted or re-flowed.  Retained uncounted as the source-local context the proof needs: lines 295--297; lines 326--332.  Nothing was omitted from any retained span, so the package carries no omission record.  No retained line contains a citation, so the wrapper carries no bibliography and no bibliography entry.  No retained line contains a figure environment, an image-inclusion command or drawing code, so no image is delivered and the package declares no asset dependency.  Editorial formatting only, in addition: an AMS \texttt{amsart} preamble that restates the article's own declarations and macros used by the retained lines, namely the environments \texttt{proposition}, \texttt{lemma} on separate counters, together with the standard packages those lines need.  The retained environments print in this document as Proposition 1, Lemma 1 in order of appearance, whereas the article prints the same items with the numbering of its own full document.  The complete native source of the article as distributed in the arXiv e-print for this exact version is bundled unchanged as \texttt{sources/200581/source0.tex}.
\begin{proposition}\label{prg}
 Any parabolic soliton to the ICSF is a graph on the $y$-axis.
\end{proposition}

\begin{equation}
\label{para2}
\begin{cases}
y' = \sin\theta \\
\displaystyle \theta' =  \frac{1}{\sin\theta} - \frac{\cos\theta}{y},
\end{cases}
\end{equation}

\begin{lemma}\label{le56}
    Let $\alpha_\gamma = (y, \theta)$ be an orbit of the system \eqref{para2} with the initial condition $\alpha_\gamma(0) = (y_0, \pi/2)$. If $\alpha_\gamma$ is of Type I, then $\lim_{s \to s_m} \alpha_\gamma(s) = (0, \pi/2)$.
\end{lemma}

\begin{proof}
    Since $\alpha_\gamma$ is a Type I orbit, and by Lemma {lem-unique-intersection}, there exists a unique $\hat{s} < 0$ such that $\alpha_\gamma(\hat{s}) \in \Gamma$. Denote $\hat{y} := y(\hat{s})$ and $\hat{\theta} := \theta(\hat{s})$. Moreover, we know that  $\alpha_\gamma$ is contained inside the region determined by $\Gamma$ for all $y \in (0, \hat{y})$. Then $\lim_{s \to s_m} \alpha_\gamma(s) = (0, \theta_\ast)$, where $\theta_\ast \in (\hat{\theta}, \pi / 2] \subseteq (0, \pi/2]$.

    Suppose that $\theta_\ast \neq \pi / 2$. By Proposition \ref{prg}, we  can consider $\theta$ as a function of $y$, i.e. $\theta = \theta(y)$ for $y \in (0, \hat{y})$. Since $\alpha_\gamma$ is contained inside $\Gamma$, we know that $\theta(y)$ is a decreasing function on the interval $(0, \hat{y})$. Therefore, we have $\theta(y) \in (\hat{\theta}, \theta_\ast)$ for all $y \in (0,\hat{y})$.

    Since the functions $\theta \mapsto \frac{1}{\sin^2 \theta}$ and $\theta \mapsto \cot\theta$ are positive and continuous on $(0, \pi/2)$, they are in particular bounded on $(0,\hat{y})$, since for $y \in (0, \hat{y})$, we have $\theta(y) \in (\hat{\theta}, \theta_\ast)$ and $\theta_\ast \neq \pi / 2$. So, there are constants $A,B > 0$ such that
    \begin{equation*}
       	\begin{split}
	    \frac{1}{\sin^2 \theta(y)} &\leq A, \quad  y \in (0, \hat{y}) \\
	    \cot(\theta(y)) &\geq B, \quad  y \in (0, \hat{y}).
       	\end{split}
    \end{equation*}

    In particular, for all $y \in (0, \hat{y})$ we have
    $$\frac{d\theta}{dy} = \frac{1}{\sin^2\theta} - \frac{\cot\theta}{y} \leq A - \frac{B}{y}.$$

    Integrating from $y$ to $\hat{y}$, $y < \hat{y}$, we get
    $$\theta(y) - \hat{\theta} \leq A(\hat{y} - y) -B\log(\hat{y}) + B\log(y).$$

    Letting $y \to 0,$ we get $$\theta_\ast - \hat{\theta} \leq - \infty,$$ which is an obvious contradiction. Hence, $\theta_\ast = \pi/2$ and so $$\lim_{s \to s_m} \theta(s) =\frac{ \pi}{2}.$$
\end{proof}

\end{document}
